%%%%%%%%%%%%%%%%%%%%%%%%%%%%%%%%%%%%%%%%%%%%%%%%%%%%%%%%%%%%%%%%%%%%%%%%
%%%%%%%%%%%%%%%%%%%%%% Simple LaTeX CV Template %%%%%%%%%%%%%%%%%%%%%%%%
%%%%%%%%%%%%%%%%%%%%%%%%%%%%%%%%%%%%%%%%%%%%%%%%%%%%%%%%%%%%%%%%%%%%%%%%

%%%%%%%%%%%%%%%%%%%%%%%%%%%%%%%%%%%%%%%%%%%%%%%%%%%%%%%%%%%%%%%%%%%%%%%%
%% NOTE: If you find that it says                                     %%
%%                                                                    %%
%%                           1 of ??                                  %%
%%                                                                    %%
%% at the bottom of your first page, this means that the AUX file     %%
%% was not available when you ran LaTeX on this source. Simply RERUN  %%
%% LaTeX to get the ``??'' replaced with the number of the last page  %%
%% of the document. The AUX file will be generated on the first run   %%
%% of LaTeX and used on the second run to fill in all of the          %%
%% references.                                                        %%
%%%%%%%%%%%%%%%%%%%%%%%%%%%%%%%%%%%%%%%%%%%%%%%%%%%%%%%%%%%%%%%%%%%%%%%%

%%%%%%%%%%%%%%%%%%%%%%%%%%%% Document Setup %%%%%%%%%%%%%%%%%%%%%%%%%%%%

% Don't like 10pt? Try 11pt or 12pt
\documentclass[11pt]{article}

% This is a helpful package that puts math inside length specifications
\usepackage{calc}

% Simpler bibsection for CV sections
% (thanks to natbib for inspiration)


% Layout: Puts the section titles on left side of page
\reversemarginpar

%
%         PAPER SIZE, PAGE NUMBER, AND DOCUMENT LAYOUT NOTES:
%
% The next \usepackage line changes the layout for CV style section
% headings as marginal notes. It also sets up the paper size as either
% letter or A4. By default, letter was used. If A4 paper is desired,
% comment out the letterpaper lines and uncomment the a4paper lines.
%
% As you can see, the margin widths and section title widths can be
% easily adjusted.
%
% ALSO: Notice that the includefoot option can be commented OUT in order
% to put the PAGE NUMBER *IN* the bottom margin. This will make the
% effective text area larger.
%
% IF YOU WISH TO REMOVE THE ``of LASTPAGE'' next to each page number,
% see the note about the +LP and -LP lines below. Comment out the +LP
% and uncomment the -LP.
%
% IF YOU WISH TO REMOVE PAGE NUMBERS, be sure that the includefoot line
% is uncommented and ALSO uncomment the \pagestyle{empty} a few lines
% below.
%

%% Use these lines for letter-sized paper
\usepackage[paper=letterpaper,
            %includefoot, % Uncomment to put page number above margin
            marginparwidth=1.2in,     % Length of section titles
            marginparsep=.05in,       % Space between titles and text
            margin=1in,               % 1 inch margins
            includemp]{geometry}

%% Use these lines for A4-sized paper
%\usepackage[paper=a4paper,
%            %includefoot, % Uncomment to put page number above margin
%            marginparwidth=30.5mm,    % Length of section titles
%            marginparsep=1.5mm,       % Space between titles and text
%            margin=25mm,              % 25mm margins
%            includemp]{geometry}

%% More layout: Get rid of indenting throughout entire document
\setlength{\parindent}{0in}

%% Set line spacing
\linespread{1.2}

%% This gives us fun enumeration environments. compactitem will be nice.
\usepackage{paralist}
%% Reference the last page in the page number
%
% NOTE: comment the +LP line and uncomment the -LP line to have page
%       numbers without the ``of ##'' last page reference)
%
% NOTE: uncomment the \pagestyle{empty} line to get rid of all page
%       numbers (make sure includefoot is commented out above)
%
\usepackage{fancyhdr,lastpage}
\pagestyle{fancy}
\pagestyle{empty}      % Uncomment this to get rid of page numbers
\fancyhf{}\renewcommand{\headrulewidth}{0pt}
\fancyfootoffset{\marginparsep+\marginparwidth}
\newlength{\footpageshift}
\setlength{\footpageshift}
          {0.5\textwidth+0.5\marginparsep+0.5\marginparwidth-2in}
\lfoot{\hspace{\footpageshift}%
       \parbox{4in}{\, \hfill %
                    \arabic{page} of \protect\pageref*{LastPage} % +LP
%                    \arabic{page}                               % -LP
                    \hfill \,}}

% Finally, give us PDF bookmarks
\usepackage{color,hyperref}
\definecolor{darkblue}{rgb}{0.0,0.0,0.3}
\hypersetup{colorlinks,breaklinks,
            linkcolor=darkblue,urlcolor=darkblue,
            anchorcolor=darkblue,citecolor=darkblue}

%%%%%%%%%%%%%%%%%%%%%%%% End Document Setup %%%%%%%%%%%%%%%%%%%%%%%%%%%%


%%%%%%%%%%%%%%%%%%%%%%%%%%% Helper Commands %%%%%%%%%%%%%%%%%%%%%%%%%%%%

% The title (name) with a horizontal rule under it
%
% Usage: \makeheading{name}
%
% Place at top of document. It should be the first thing.
\newcommand{\makeheading}[1]%
        {\hspace*{-\marginparsep minus \marginparwidth}%
         \begin{minipage}[t]{\textwidth+\marginparwidth+\marginparsep}%
                {\Large \bfseries #1}\\[-0.15\baselineskip]%
                 \rule{\columnwidth}{1pt}%
         \end{minipage}}

% The section headings
%
% Usage: \section{section name}
%
% Follow this section IMMEDIATELY with the first line of the section
% text. Do not put whitespace in between. That is, do this:
%
%       \section{My Information}
%       Here is my information.
%
% and NOT this:
%
%       \section{My Information}
%
%       Here is my information.
%
% Otherwise the top of the section header will not line up with the top
% of the section. Of course, using a single comment character (%) on
% empty lines allows for the function of the first example with the
% readability of the second example.
\renewcommand{\section}[2]%
        {\pagebreak[2]\vspace{\baselineskip}%
         \phantomsection\addcontentsline{toc}{section}{#1}%
         \hspace{0in}%
         \marginpar{
         \raggedright \scshape #1}#2}

% An itemize-style list with lots of space between items
\newenvironment{outerlist}[1][\enskip\textbullet]%
        {\begin{itemize}[#1]}{\end{itemize}%
         \vspace{-0.8\baselineskip}}

% An environment IDENTICAL to outerlist that has better pre-list spacing
% when used as the first thing in a \section
\newenvironment{lonelist}[1][\enskip\textbullet]%
        {\vspace{-\baselineskip}\begin{list}{#1}{%
        \setlength{\partopsep}{0pt}%
        \setlength{\topsep}{0pt}}}
        {\end{list}\vspace{-.6\baselineskip}}

% An itemize-style list with little space between items
\newenvironment{innerlist}[1][\enskip\textbullet]%
        {\begin{compactitem}[#1]}{\end{compactitem}}

% An environment IDENTICAL to innerlist that has better pre-list spacing
% when used as the first thing in a \section
\newenvironment{loneinnerlist}[1][\enskip\textbullet]%
        {\vspace{-\baselineskip}\begin{compactitem}[#1]}
        {\end{compactitem}\vspace{-.6\baselineskip}}

% To add some paragraph space between lines.
% This also tells LaTeX to preferably break a page on one of these gaps
% if there is a needed pagebreak nearby.
\newcommand{\blankline}{\quad\pagebreak[2]}

%%%%%%%%%%%%%%%%%%%%%%%% End Helper Commands %%%%%%%%%%%%%%%%%%%%%%%%%%%

%%%%%%%%%%%%%%%%%%%%%%%%% Begin CV Document %%%%%%%%%%%%%%%%%%%%%%%%%%%%

\begin{document}
\makeheading{Tiffany Bao}

\section{Contact Information}
%
% NOTE: Mind where the & separators and \\ breaks are in the following
%       table.
%
% ALSO: \rcollength is the width of the right column of the table
%       (adjust it to your liking; default is 1.85in).
%
%\newlength{\rcollength}\setlength{\rcollength}{1in}%
%
%\begin{tabular}[t]{@{}p{\textwidth-\rcollength}p{\rcollength}}
699 South Mill Avenue \\
Arizona State University  \\
Tempe, AZ 85281 \\
Email: \href{mailto:tbao@asu.edu}{\tt{tbao@asu.edu}}
%Website: \url{https://users.ece.cmu.edu/~youzhib}
%\end{tabular}

\section{Professional Appointments}
\textbf{Associate Professor, Arizona State University}
\hfill{2025 August -- present}\\
\textbf{Assistant Professor, Arizona State University}
\hfill{2018 -- present}\\

%\section{Research Interests}
%
%Binary Analysis, Software Security and Security Economics.

%\blankline

%%%%%%%%%%%%%%%%%%%
%%%% Education %%%%
%%%%%%%%%%%%%%%%%%%

\section{Education}
%
\textbf{Doctor of Philosophy, Carnegie Mellon University}
\hfill{2012 -- 2018}
%\begin{outerlist}
%\item[] Department of Electrical and Computer
    %Engineering \\
%\end{outerlist}

%\blankline

%\textbf{Bachelor of Science, Peking University}
%\hfill{2008 -- 2012}
%\begin{outerlist}
%\item[] School of Electrical Engineering \& Computer Science \\
    %Department of Computer Science
%\end{outerlist}

\section{Publications}
Teamwork Makes the Defense Work: Comprehensive Vulnerability Defense Resource Allocation. \\
Siyu Liu, Rida Bazzi, Fei Fang, Tiffany Bao. \\
In \emph{Proceedings of the 24th International Conference on Autonomous Agents and Multiagent Systems (AAMAS '25)}.
\\\\
``It's almost like Frankenstein'': Investigating the Complexities of Scientific Collaboration and Privilege Management within Research Computing Infrastructures. \\
Souradip Nath, Ananta Soneji, Jaejong Baek, Tiffany Bao, Adam Doup\'{e}, Carlos Rubio-Medrano, Gail-Joon Ahn. \\
In \emph{Proceedings of the 46th IEEE Symposium on Security and Privacy (Oakland '25)}.
\\\\
ScamNet: Toward Explainable Large Language Model-Based Fraudulent Shopping Website Detection. \\
Marzieh Bitaab, Alireza Karimi, Zhuoer Lyu, Ahmadreza Mosallanezhad, Adam Oest, Ruoyu Wang, Tiffany Bao, Yan Shoshitaishvili, Adam Doup\'{e}. \\
In \emph{Proceedings of the AAAI Conference on Artificial Intelligence (AAAI '25)}.
\\\\
System Register Hijacking: Compromising Kernel Integrity By Turning System Registers Against the System. \\
Jennifer Miller, Manas Ghandat, Kyle Zeng, Hongkai Chen, Abdelouahab Habs Benchikh, Tiffany Bao, Ruoyu Wang, Adam Doup\'{e}, Yan Shoshitaishvili. \\
In \emph{Proceedings of the 34th USENIX Security Symposium (USENIX '25)}.
\\\\
SCAMMAGNIFIER: Piercing the Veil of Fraudulent Shopping Website Campaigns. \\
Marzieh Bitaab, Alireza Karimi, Zhuoer Lyu, Adam Oest, Dhruv Kuchhal, Muhammad Saad, Gail-Joon Ahn, Ruoyu Wang, Tiffany Bao, Yan Shoshitaishvili, Adam Doup\'{e}. \\
In \emph{Proceedings of the Network and Distributed System Security Symposium (NDSS '25)}.
\\\\
Fuzz to the Future: Uncovering Occluded Future Vulnerabilities via Robust Fuzzing. \\
Arvind S Raj, Wil Gibbs, Fangzhou Dong, Jayakrishna Menon Vadayath, Michael Tompkins, Steven Wirsz, Yibo Liu, Zhenghao Hu, Chang Zhu, Gokulkrishna Praveen Menon, Brendan Dolan-Gavitt, Adam Doup\'{e}, Ruoyu Wang, Yan Shoshitaishvili, Tiffany Bao. \\
In \emph{Proceedings of the 2024 ACM SIGSAC Conference on Computer and Communications Security (CCS '24)}.
\\\\
TYGR: Type Inference on Stripped Binaries using Graph Neural Networks. \\
Chang Zhu, Ziyang Li, Anton Xue, Ati Priya Bajaj, Wil Gibbs, Yibo Liu, Rajeev Alur, Tiffany Bao, Hanjun Dai, Adam Doup\'{e}, Mayur Naik, Yan Shoshitaishvili, Ruoyu Wang, Aravind Machiry. \\
In \emph{Proceedings of the 33rd USENIX Security Symposium (USENIX '24)}.
\\\\
Take a Step Further: Understanding Page Spray in Linux Kernel Exploitation. \\
Ziyi Guo, Dang K Le, Zhenpeng Lin, Kyle Zeng, Ruoyu Wang, Tiffany Bao, Yan Shoshitaishvili, Adam Doup\'{e}, Xinyu Xing. \\
In \emph{Proceedings of the 33rd USENIX Security Symposium (USENIX '24)}.
\\\\
Ahoy SAILR! There is No Need to DREAM of C: A Compiler-Aware Structuring Algorithm for Binary Decompilation. \\
Zion Leonahenahe Basque, Ati Priya Bajaj, Wil Gibbs, Jude O'Kain, Derron Miao, Tiffany Bao, Adam Doup\'{e}, Yan Shoshitaishvili, Ruoyu Wang. \\
In \emph{Proceedings of the 33rd USENIX Security Symposium (USENIX '24)}.
\\\\
AirTaint: Making Dynamic Taint Analysis Faster and Easier. \\
Qian Sang, Yanhao Wang, Yuwei Liu, Xiangkun Jia, Tiffany Bao, Purui Su. \\
In \emph{Proceedings of the IEEE Symposium on Security and Privacy (Oakland '24)}.
\\\\
“Watching over the shoulder of a professional”: Why hackers make mistakes and how they fix them. \\
Irina Ford, Ananta Soneji, Faris Bugra Kokulu, Jayakrishna Vadayath, Zion Leonahenahe Basque, Gaurav Vipat, Adam Doup\'{e}, Ruoyu Wang, Gail-Joon Ahn, Tiffany Bao, Yan Shoshitaishvili. \\
In \emph{Proceedings of the IEEE Symposium on Security and Privacy (Oakland '24)}.
\\\\
RetSpill: Igniting User-Controlled Data to Burn Away Linux Kernel Protections. \\
Kyle Zeng, Zhenpeng Lin, Kangjie Lu, Xinyu Xing, Ruoyu Wang, Adam Doup\'{e}, Yan Shoshitaishvili, Tiffany Bao. \\
In \emph{Proceedings of the 2023 ACM SIGSAC Conference on Computer and Communications Security (CCS '23)}.
\\\\
Greenhouse: Single-Service Rehosting of Linux-Based Firmware Binaries in User-Space Emulation. \\
Hui Jun Tay, Kyle Zeng, Jayakrishna Menon Vadayath, Arvind S Raj, Audrey Dutcher, Tejesh Reddy, Wil Gibbs, Zion Leonahenahe Basque, Fangzhou Dong, Zack Smith, Adam Doup\'{e}, Tiffany Bao, Yan Shoshitaishvili, Ruoyu Wang. \\
In \emph{the 32nd USENIX Security Symposium (USENIX '23)}.
\\\\
BEYOND PHISH: Toward Detecting Fraudulent e-Commerce Websites at Scale. \\
Marzieh Bitaab, Haehyun Cho, Adam Oest, Zhuoer Lyu, Wei Wang, Jorij Abraham, Ruoyu Wang, Tiffany Bao, Yan Shoshitaishvili, Adam Doup\'{e}. \\
In \emph{Proceedings of the IEEE Symposium on Security and Privacy (Oakland '23)}.
\\\\
Toss a Fault to Your Witcher: Applying Grey-box Coverage-guided Mutational Fuzzing to Detect SQL and Command Injection Vulnerabilities. \\
Erik Trickel, Fabio Pagani, Chang Zhu, Lukas Dresel, Giovanni Vigna, Christopher Kruegel, Ruoyu Wang, Tiffany Bao, Yan Shoshitaishvili, Adam Doup\'{e}. \\
In \emph{Proceedings of the IEEE Symposium on Security and Privacy (Oakland '23)}.
\\\\
Cyber Deception: Techniques, Strategies, and Human Aspects.\\
Tiffany Bao, Milind Tambe, Cliff Wang.\\
Springer Nature.
\\\\
Mitigating Threats Emerging from the Interaction between SDN Apps and SDN (Configuration) Datastore. \\
Sana Habib, Tiffany Bao, Yan Shoshitaishvili, Adam Doup\'{e}. \\
Proceedings of the 2022 on Cloud Computing Security Workshop.
\\\\
Playing for K(H)eaps: Understanding and Improving Linux Kernel Exploit Reliability. \\
Kyle Zeng, Yueqi Chen, Haehyun Cho, Xinyu Xing, Adam Doup\'{e}, Yan Shoshitaishvili, Tiffany Bao. \\
In \emph{the 31st USENIX Security Symposium (USENIX '22)}.
\\\\
Arbiter: Bridging the Static and Dynamic Divide in Vulnerability Discovery on Binary Programs. \\
Jayakrishna Vadayath, Moritz Eckert, Kyle Zeng, Nicolaas Weideman, Gokulkrishna Praveen Menon, Yanick Fratantonio, Davide Balzarotti, Adam Doup\'{e}, Tiffany Bao, Ruoyu Wang, Christophe Hauser, Yan Shoshitaishvili. \\
In \emph{the 31st USENIX Security Symposium (USENIX '22)}.
\\\\
Expected Exploitability: Predicting the Development of Functional Vulnerability Exploits. \\
Octavian Suciu, Connor Nelson, Zhuoer Lyu, and Tiffany Bao, Tudor Dumitra\c{s}. \\
In \emph{the 31st USENIX Security Symposium (USENIX '22)}.
\\\\
"Flawed, but like democracy we don't have a better system": The Experts' Insights on the Peer Review Process of Evaluating Security Papers. \\
Ananta Soneji, Faris Bugra Kokulu, Carlos Rubio-Medrano, Tiffany Bao, Ruoyu Wang, Yan Shoshitaishvili, Adam Doup\'{e}.\\
In \emph{Proceedings of the IEEE Symposium on Security and Privacy (Oakland '22)}.
\\\\
ViK: practical mitigation of temporal memory safety violations through object ID inspection. \\
Haehyun Cho, Jinbum Park, Adam Oest, Tiffany Bao, Ruoyu Wang, Yan Shoshitaishvili, Adam Doup\'{e}, Gail-Joon Ahn. \\
Proceedings of the 27th ACM International Conference on Architectural Support for Programming Languages and Operating Systems (ASPLOS '22).
\\\\
Favocado: Fuzzing Binding Code of JavaScript Engines Using Semantically Correct Test Cases. \\
Sung Ta Dinh, Haehyun Cho, Kyle Martin, Adam Oest, Yihui Zeng, Alexandros Kapravelos, Tiffany Bao, Ruoyu Wang, Yan Shoshitaishvili, Adam Doup\'{e}, Gail-Joon Ahn. \\
In \emph{Network and Distributed System Security Symposium (NDSS '21)}.
\\\\
Having Your Cake and Eating It: An Analysis of Concession-Abuse-as-a-Service \\
Zhibo Sun, Adam Oest, Penghui Zhang, Carlos Rubio-Medrano, Tiffany Bao, Ruoyu Wang, Ziming Zhao, Yan Shoshitaishvili, Adam Doup\'{e}, Gail-Joon Ahn. \\
In \emph{the 30th USENIX Security Symposium (USENIX '21)}.
\\\\
CrawlPhish: Large-scale Analysis of Client-side Cloaking Techniques in Phishing. \\
Penghui Zhang, Adam Oest, Haehyun Cho, Zhibo Sun, RC Johnson, Brad Wardman, Shaown Sarker, Alexandros Kapravelos, 
Tiffany Bao, Ruoyu Wang, Yan Shoshitaishvili, Adam Doup\'{e}, Gail-Joon Ahn. \\
In \emph{Proceedings of the IEEE Symposium on Security and Privacy (Oakland '21)}.
\\\\
SyML: Guiding Symbolic Execution Toward Vulnerable States Through Pattern Learning. \\
Nicola Ruaro, Kyle Zeng, Lukas Dresel, Mario Polino, Tiffany Bao, Andrea Continella, Stefano Zanero, Christopher Kruegel, Giovanni Vigna. \\
In \emph{the 24th International Symposium on Research in Attacks, Intrusions and Defenses (RAID '21)}.
\\\\
SoK: Everything You Ever Wanted to Know About Bitcoin Mixers (But Were Afraid to Ask) \\
Jaswant Pakki, Yan Shoshitaishvili, Ruoyu Wang, Tiffany Bao, Adam Doup\'{e}. \\
In \emph{Financial Cryptography and Data Security (FC '21)}.
\\\\
MuTent: Dynamic Android Intent Protection with Ownership-Based Key Distribution and Security Contracts. \\
Pradeep Kumar Duraisamy Soundrapandian, Jaejong Baek, Tiffany Bao, Yan Shoshitaishvili, Adam Doup\'{e}, Ruoyu Wang, Gail-Joon Ahn. \\
In \emph{Hawaii International Conference on System Sciences (HICCS '21)}. \\
\textbf{Best Minitrack Paper Award}.
\\\\
HoneyPLC: A Next-Generation Honeypot for Industrial Control Systems \\
Efren L\'{o}pez-Morale, Carlos Rubio-Medrano, Adam Doup\'{e}, Yan Shoshitaishvili, Ruoyu Wang, Tiffany Bao, Gail-Joon Ahn.\\
In \emph{Proceedings of the 2020 ACM SIGSAC Conference on Computer and Communications Security (CCS '20)}.
\\\\
Exploiting Uses of Uninitialized Stack Variables in Linux Kernels to Leak Kernel Pointers. \\
Haehyun Cho, Jinbum Park, Joonwon Kang, Tiffany Bao, Ruoyu Wang, Yan Shoshitaishvili, Adam Doup\'{e}, Gail-Joon Ahn. \\
In \emph{the 14th USENIX Workshop on Offensive Technologies (WOOT '20)}.
\\\\
Not All Coverage Measurements Are Equal: Fuzzing by Coverage Accounting for Input Prioritization. \\
Yanhao Wang, Xiangkun Jia, Yuwei Liu, Tiffany Bao, Dinghao Wu, and Purui Su. \\
In the \emph{Network and Distributed System Security Symposium (NDSS '20)}.
\\\\
Scam Pandemic: How Attackers Exploit Public Fear through Phishing. \\
Marzieh Bitaab, Haehyun Cho, Adam Oest, Penghui Zhang, Zhibo Sun, Rana Pourmohamad, Doowon Kim, Tiffany Bao, Ruoyu Wang, Yan Shoshitaishvili, Adam Doup\'{e}, Gail-Joon Ahn.
In \emph{Symposium on Electronic Crime Research (ECrime '20)}.
\\\\
%
Cyber Autonomy in Software Security: Techniques and Tactics (Book Chapter). \\
Tiffany Bao and Yan Shoshitaishvili. \\
In \emph{Game Theory and Machine Learning for Cyber Security}.
\\\\
Matched and Mismatched SOCs: A Qualitative Study on Security Operations Center Issues. \\
Faris Bugra Kokulu, Ananta Soneji, Tiffany Bao, Yan Shoshitaishvili, Ziming Zhao, Adam Doup\'{e}, and Gail-Joon Ahn. \\
In \emph{Proceedings of the 2019 ACM SIGSAC Conference on Computer and Communications Security (CCS '19)}.
\\\\
%
Understanding and Predicting Private Interactions in Underground Forums. \\
Zhibo Sun, Carlos E. Rubio-Medrano, Ziming Zhao, Tiffany Bao, Adam Doup\'{e}, and Gail-Joon Ahn.\\
In \emph{the 9th ACM Conference on Data and Application Security and Privacy (CODASPY '19)}.
\\\\
%%
%Autonomous Computer Security Game: Techniques, Strategy and Investigation. \\
%\textbf{Tiffany Bao} \\
%Doctor of Philosophy Thesis.
%
Your Exploit is Mine: Automatic Shellcode Transplant for
Remote Exploits.\\
Tiffany Bao, Ruoyu Wang, Yan Shoshitaishvili, and David
Brumley.\\
In \emph{Proceedings of the 38th IEEE Symposium on
Security and Privacy (Oakland '17)}.
\\\\
%
How Shall We Play a Game: A Game-theoretical Model for
Cyber-warfare.\\
Tiffany Bao, Yan Shoshitaishvili, Ruoyu Wang, Christopher Kruegel,
Giovanni Vigna and David Brumley. \\
In \emph{Proceedings of the 30th IEEE Computer Security Foundations
(CSF '17)}. \\
\textbf{National Security Agency's Annual Scientific Paper Award.}
\\\\
%
Security is a Game. \\
Tiffany Bao. \\
In \emph{2017 USENIX Summit on Hot Topics in Security (HotSec '17)}. \\
\textbf{Awarded Talk.}
\\\\
%
A Game-theoretical Model for Cyber-warfare Games. (Invited
Poster)\\
Tiffany Bao, Yan Shoshitaishvili, Ruoyu Wang and David
Brumley. \\
In \emph{the 8th Workshop on Computational Cybersecurity in
Compromised Environments (C3E '16)}.
\\\\
%
ByteWeight: Learning to Recognize Functions in Binary Code.
\\
Tiffany Bao, Jonathan Burket, Maverick Woo, Rafael Turner,
and David Brumley.  \\
In \emph{Proceedings of the 23rd USENIX Security Symposium
(USENIX '14)}.
\\\\
%
Type-based Dynamic Taint Analysis Technology. \\
Libo Chen, Jianwei Zhuge, Fan Tian, Tiffany Bao, and Xun Lu.
In \emph{Tsinghua Science and Technology Journal}, 2012.
\\\\
%
Research of Technology for Type-based Dynamic Taint Analysis. \\
Libo Chen, Jianwei Zhuge, Fan Tian, Tiffany Bao, and Xun Lu.
In \emph{Proceedings of the 5th Conference of Vulnerability Analysis
and Risk Assessment (VARA '12)}. \\
\textbf{Outstanding Paper Award.}

%\\\\

%Efficient Crash Triage with Black-box Dependency Inference. \\
%\textbf{Tiffany Bao}, Thanassis Avgerinos, Gustavo Grieco and David Brumley. \\
%Submitted for publication.
%\\\\
%Primus: Memory Check with Micro Execution. \\
%Ivan Gotovchits, \textbf{Tiffany Bao} and David Brumley. \\
%Submitted for publication.
%\\\\
%Does IT Make us Secure: A Qualitative Evaluation for Binary Analysis. \\
%\textbf{Tiffany Bao} and David Brumley. \\
%To be submitted.



\section{Funded Projects}
CAREER: Achieving Autonomous Symbolic Execution through Learning from Humans \hfill{2025-2030} \\
National Science Foundation (NSF). \\
Awarded funding: \$500,000 
\\\\
SENPAI: Strategic Exploration, Navigation, and Patching of Abstracted Integrations for Cyber-Physical Systems	\hfill{2024-2027}
DOD-DARPA: Information Innovation Office (DARPA MTO). \\
Awarded funding: 	\$7,960,690 
\\\\
EURYALE: Combating Emergent Execution with a GLANCE \hfill{2025-2026} \\
DOD-DARPA: Information Innovation Office (IIO). \\
Awarded funding: \$1,523,727 \\
\\
Dynamic and Personalized Deception for Proactive Cyber Defense \hfill{2024-2027} \\
DOD-USAF-AFRL: Air Force Office of Scientific Research (AFOSR). \\
Awarded funding: \$599,968 \\
\\
American Cybersecurity Education (ACE) Institute \hfill{2024-2026} \\
DOD: Defense Advanced Research Projects Agency (DARPA). \\
Awarded funding: \$4,462,558 \\
\\
RxCRS: Reliable and eXplainable Cyber Reasoning System for Digital Health Security \hfill{2023-2025} \\
HHS-NIH: Advanced Research Projects Agency for Health (ARPA-H). \\
Awarded funding: \$6,564,290 \\
\\
Assessing Cyber Training by Cyber Operation Mistakes \hfill{2023-2026} \\
DOD-NAVY: Office of Naval Research (ONR). \\
Awarded funding: \$900,135 \\
\\
SaTC: CORE: Medium: Symbolizing Viability: Paving the Road to Practical Symbolic Execution \hfill{2023-2027} \\
National Science Foundation (NSF). \\
Awarded funding: \$1,199,992 \\
\\
NSF-SFS: Arizona Cyber Defense Scholarship \hfill{2017-2026} \\
National Science Foundation (NSF). \\
Awarded funding: \$4,297,774 \\
Cybersecurity Meta-analysis for USAID Beneficiaries. \hfill{2021-2022} \\
US Agency for International Development (USAID). \\
Awarded funding: \$59,655
\\\\

Predicting and Assessing Attacks for Open Source Projects. \hfill{2021-2022} \\
DOD-DARPA: Information Innovation Office (DARPA IIO). \\
Awarded funding: \$487,720
\\\\

SaTC: CORE: Medium: Self-Adaptive Cyber Risk Management  \hfill{2020 - 2023} \\ 
via Machine to Machine Economy supported by blockchain and \\
smart contracts technology.\\
National Science Foundation (NSF). \\
Awarded funding: \$750,000 
\\\\


VOLT: A Viscous, Orchestrated Lifting and Translation Framework	\hfill{2020-2023} \\
DOD-DARPA: Information Innovation Office (DARPA IIO). \\
Awarded funding: \$5,045,355
\\\\


Human-Assisted Cyber Reasoning Systems and Oppositional \hfill{2019-2023} \\
Human Factors. \\
DOD: National Security Agency (NSA). \\
Awarded funding: \$5,598,441
\\\\

CHECRS: Cognitive Human Enhancements for Cyber Reasoning \hfill{2018-2022} \\
Systems.  \\
DOD-DARPA: Information Innovation Office (DARPA IIO). \\
Awarded funding: \$11,730,556
\\\\

Realizing Cyber Inception: Towards a Science of Personalized \hfill{2017-2023} \\
Deception for Cyber Defense. \\
DOD-ARMY-ARL: Army Research Office (ARO). \\
Awarded funding: \$598,149


%\input{talks}


\section{Professional Activities}
%
\textbf{Chair Member} \\
Research in Attacks, Intrusions and Defenses. Co-Chair. \hfill{2025, 2026} \\
Network and Distributed System Security Symposium Poster Session. \hfill{2023, 2024} \\
USENIX Security, Associate chair member. \hfill{2024} \\
IEEE Security and Privacy, Vice chair member. \hfill{2024} \\
\\
%
\textbf{Program Committee Member} \\
Network and Distributed System Security Symposium. \hfill{2019 - 2023, 2025} \\
ACM SIGSAC Conference on Computer and Communications Security. \hfill{2020, 2022-2024} \\
IEEE Security and Privacy. \hfill{2022, 2023, 2024} \\
USENIX Security. \hfill{2022, 2024, 2025} \\
International Symposium on Research in Attacks, Intrusions and Defenses. \hfill{2020}
ACM Workshop on Automotive and Aerial Vehicle Security. \hfill{2020} \\
(Affiliated with the ACM Conference on Data and Application Security and Privacy.) \\
Network and Distributed System Security Symposium. \hfill{2019} \\
Annual Computer Security Applications Conference. \hfill{2019} \\
Workshop on Binary Analysis Research. \hfill{2018, 2019, 2020} \\
(Affiliated with the Network and Distributed System Security Symposium.) \\
%
\\
\textbf{Journal Editorial Board} \\
The Journal of Systems Research, Security Track. \\
\\
%
\textbf{Journal Reviewer} \\
ACM Transactions on information and System Security. \\
Communications of the ACM. \\
IEEE Transactions on Software Engineering. \\
IEEE Transactions on Knowledge and Data Engineering. \\
\\
\textbf{CTF Game Organizer} \\
Member of the DEF CON CTF hosting team. \hfill{2019 - 2021}
%
\\\\
%\textbf{Student Committee Member} \\
%IEEE Symposium on Security and Privacy. \hfill{2016}
%\\\\
%\textbf{External Reviewer} \\
%The Network and Distributed System Security Symposium.      \hfill{2014}
%\\
%The ACM Conference on Computer and Communications Security.  \hfill{2013}
%\\\\
\textbf{Mentor} \\
Google Summer of Code.  \hfill{2013}
%\\\\
%\textbf{Contributor} \\
%\emph{Introduction to Algorithm and Olympiad in Informatics.} \hfill{2011}\\
%ISBN:9787302291077
%\\\vspace{-1em}\\
%\emph{Network Hacking and Defense.}\hfill{2010}\\
%ISBN: 9787121138027

\section{Selected Awards}
%
% \href{http://english.pku.edu.cn/}{Peking University}
% \begin{innerlist}
% \item
% \item Distinguished Student Award, 2008 \& 2009
% \item Social Worker Award, 2010
% \end{innerlist}
IEEE Security and Privacy Best Student Paper Award.
\hfill{2020}
\\\\
National Security Agency's Annual Scientific Paper Award.
\hfill{2018}
\\\\
Awarded Talk at 2017 USENIX Summit on Hot Topics in Security.
\hfill{2017}
\\\\
Carnegie Mellon University Presidential Fellowship.
\hfill{2017}
%%\\\\
%%Second Place at EFF Capture the Flag @ Enigma.
%%\hfill{2016}
%\\\\
%Frank J. Marshall Graduate Fellowship.
%\hfill{2012}
%\\\\
%Zongjie Huo Fellowship.
%\hfill{2008, 2009}


%\section{Teaching Experience}
%CSE 365: Introduction to Information Assurance. \hfill{2021, 2022} \\
%CSE 545: Software Security. \hfill{2020} \\
%CSE 591: Computer Security: Techniques and Tactics. \hfill{2019} \\
%\textbf{Head Teaching Assistant.} \\
%Network Security. (CMU 18-731, for Graduate) \hfill{2015}\\
%\textbf{Head Teaching Assistant.} \\
%Introduction to Computer Security. (CMU 18-487, for Undergraduate) \hfill{2016}
%\\\\
%\textbf{Instructor for Computing Olympiad.} \hfill{2008}
%%%%%%%%%%%%%%%%%%%%%%%%%%%%%%%%%
%%%% Professional Experience %%%%
%%%%%%%%%%%%%%%%%%%%%%%%%%%%%%%%%

%\section{Professional Experience}
%Research Internship, University of California Santa Barbara.  \hfill{2016}
%\\\\
%Research Internship, Tsinghua University.  \hfill{2012}
%\\\\
%Member, the Honeynet Project.  \hfill{2011-2013}
%\\\\
%Research Assistant, Peking University.  \hfill{2010-2012}

%\section{References}
%\textbf{Professor David Brumley} \\
%Carnegie Mellon University \\
%(412) 268-3851 \\
%dbrumley@cmu.edu
%\\\\
%\textbf{Professor Giovanni Vigna} \\
%University of California, Santa Barbara \\
%(805)-637-6382 \\
%vigna@cs.ucsb.edu
%\\\\
%\textbf{Professor Hui Zhang} \\
%Carnegie Mellon University \\
%(412) 443-8196 \\
%hzhang@cs.cmu.edu
%\\\\
%\textbf{Professor Wenke Lee} \\
%Georgia Institute of Technology \\
%(404) 385-2879 \\
%wenke.lee@gmail.com
\end{document}

%%%%%%%%%%%%%%%%%%%%%%%%%% End CV Document %%%%%%%%%%%%%%%%%%%%%%%%%%%%%
